\section{Kapitel 1: Sensoren}

\begin{multicols*}{2}
	
\begin{itemize}
	
	\item Messung von linearer Position bzw. Drehwinkel
	\subitem Schiebe- bzw. Linearpotentiometer
	\[ U_a = U_B \cdot \frac{\Delta R}{R} = U_b \cdot \frac{l_x}{l} \]
	\subitem Drehpotentiometer (0° - ca 350°)
	\[ U_s = U_0 \cdot \frac{\Delta R}{R} = U_0 \cdot \frac{\alpha}{\alpha_0} \]
	
	\vspace{1cm}
	
	\item Drehzahlmessung mittels Inkrementalgebern
	\subitem Zählung der Pulse N pro Zeit T
	\[ \omega = \frac{\Delta\varphi \cdot N}{N} \]
	\[ \Delta\omega = \frac{\Delta\varphi}{T}\]
	\subitem Messung der Zeitdifferenz T zwischen zwei aufeinanderfolgenden Pulsen
	\[ \omega = \frac{\Delta\varphi}{T} \]
	\[ \Delta\omega = \frac{\Delta\varphi}{T^2} \cdot \Delta T \]

	\vspace{1cm}

	\item Generierung Gray-Code\\
	\\
	X1: Dualzahl im Binärcode\\
	X2: X1 >> 1 (Rechtsshift um 1)\\
	X3: X1 xor X2

	\vspace{1cm}
	
	\item Elektrische Widerstand eines Leiters $R(\theta)$ \textbf{temperaturabhängig}
	\[ R(\vartheta) = [1 + \alpha\cdot(\vartheta - \vartheta_0) + \beta\cdot(\vartheta - \vartheta_0)^2 + ...] \]
	\[ \overline{\alpha} = \frac{R(\vartheta_0) - R(\vartheta_u)}{(\vartheta_u - \vartheta_u)\cdot R(\vartheta_u)} \]
	$\vartheta_0$ = Bezugstemperatur\\
	\\
	Pt100 $\rightarrow R_0$ = 100 $\Omega$ \\
	Pt100 $\rightarrow R_0$ = 1000 $\Omega$ \\
	Temperaturkoeff. $\alpha$ = 3850 $\frac{ppm}{^\circ C}$

	\vspace{1cm}

	\item Seebeck-Effekt
	\[ U = U_B - U_A = \int_{T_1}^{T_2} S_B(T) - S_A(T) dT \]
	\[ = (S_B - S_A) \cdot (T_2 - T_1) \ \]

	\vspace{1cm}
	
	\item  NTC-Temperatursensoren
	\[ R(\vartheta) = R(\vartheta_0)\cdot e^{B \cdot \frac{1}{T} - \frac{1}{T_0}} \]
	$\vartheta$ = Temperatur in °C\\
	T = Abs. Temperatur\\
	
	\vspace{1cm}
	
	\item Wiederstand eines Leitungsstücks
	\[ R = \rho \cdot \frac{l}{A} = \rho \cdot \frac{4l}{\pi D^2} \]
	\[ \Delta R = \frac{\partial R}{\partial \rho} \cdot \Delta\rho + \frac{\partial R}{\partial D} \cdot \Delta D + \frac{\partial R}{\partial l} \cdot \Delta l \]
	\[ \frac{\Delta R}{R} = \left( \frac{\Delta \rho}{\rho\epsilon} + 2\mu + 1 \right)\cdot\epsilon = k \cdot \epsilon \]
	
	\vspace{1cm}
	
	\item Viertelbrücke
	\[U_{AB} = \frac{\Delta R}{4R + 2\Delta R}\cdot U_0 \]
	Wenn $\Delta$R << R
	\[U_{AB} \approx \frac{\Delta R}{4R}\cdot U_0 = \frac{U_0}{4}\cdot k \cdot \epsilon\]
	
	\vspace{1cm}
	
	\item Halbbrücke\\
	DMS auf \textbf{unterschiedlichen Trägern} aus gleichem Material, nur ein Träger wird belastet\\
	$\rightarrow$ filtert Temperaturstörungen heraus
	\[U_{AB} = \frac{\Delta R}{4R + 2\Delta R}\cdot U_0 \]
	Wenn $\Delta$R << R
	\[U_{AB} \approx \frac{\Delta R}{4R}\cdot U_0 = \frac{U_0}{4}\cdot k \cdot \epsilon\]
	
	DMS auf \textbf{gleichem Träger}\\
	$\rightarrow$ filtert Temperaturstörungen heraus und verdoppelt den Empfindlichkeit
	\[U_{AB} \approx \frac{\Delta R}{2R}\cdot U_0 = \frac{U_0}{2}\cdot k \cdot \epsilon\]
	
	\vspace{1cm}
	
	\item Vollbrücke\\
	$\rightarrow$ 4 DMS
	\[U_{AB} = \frac{\Delta R}{R} \cdot U_0 = U_0 \cdot k \cdot \epsilon\]
	
	\vspace{1cm}
	
	\item Kraftmessung (Poisson-Halbbrücke)
	\[ U_{AB} \approx (1 + v) \cdot \frac{\Delta R}{4R}\cdot U_0 = (1 + v) \cdot \frac{U_0}{4}\cdot k \cdot \epsilon \]
	$v = Poissonzahl = -\frac{\epsilon_y}{\epsilon_x}$
	
	\vspace{1cm}
	
	\item Piezoeffekt
	\[ Q = k_p \cdot F\]
	
	\vspace{1cm}
	
	\item Auswertungsschaltung Piezo-Kraftmesser
	\[ U_q = \frac{Q}{C_{ges}} = \frac{k_p \cdot F}{C_0 + C_L + C_e}\]
	
	\vspace{1cm}
	
	\item Ladungsverstärker
	\[u_a(t) = -\frac{Q}{C} = - \frac{k_p \cdot F}{C}\]
	
	\vspace{1cm}
	
	\item Barometrische Höhenmessung
	\[ h = \frac{T_0}{k_T}\cdot \left[ 1 - \left( \frac{p}{p_0}^{\frac{R_s \cdot k_T}{g}} \right) \right] + h_0\]
\end{itemize}

\end{multicols*}